\documentclass[12pt]{article}
\usepackage[utf8]{inputenc}
\usepackage[T1]{fontenc}
\usepackage[margin=1in]{geometry}
\usepackage{setspace}
\usepackage{xcolor}
\usepackage[colorlinks=true, urlcolor=blue!60!black, linkcolor=blue!60!black]{hyperref}
\setlength{\parindent}{2em}
\setlength{\parskip}{1.25em}
\renewcommand{\baselinestretch}{1.0}
\begin{document}
\begin{center}
{\LARGE \textit{Heart of Atlanta Motel, Inc. v. United States}, 379 U.S. 241 (1964)\par}
\vspace{0.8em}
{\large Case Brief\par}
\end{center}
\vspace{1em}

\section{Memory Jogger}
Congress may use the Commerce Clause to ban racial discrimination in private motels serving interstate travelers, even if its primary motivation is to remedy a moral wrong rather than purely economic harm (\href{https://www.courtlistener.com/opinion/106950/heart-of-atlanta-motel-inc-v-united-states/}{CourtListener}).

\section{Detailed Case Facts}
On July 2, 1964, Congress enacted the Civil Rights Act of 1964, 78 Stat.\ 241 (\href{https://supreme.justia.com/cases/federal/us/379/241/}{Justia}). Title II of the Act (\S\S~201--207, 42 U.S.C.\ \S\S~2000a--2000a-6) prohibited discrimination or segregation on the ground of race, color, religion, or national origin in places of public accommodation whose operations ``affect commerce'' or whose discrimination is supported by state action (\href{https://www.law.cornell.edu/supremecourt/text/379/241}{Cornell LII}). Under \S~201(b)(1) and \S~201(c)(1), any inn, hotel, motel, or other establishment providing lodging to transient guests---other than an owner-occupied building containing not more than five rooms for rent or hire---was declared to affect commerce \textit{per se} (\href{https://supreme.justia.com/cases/federal/us/379/241/}{Justia}).

Appellant Heart of Atlanta Motel, Inc., owned and operated by Moreton Rolleston, Jr., was a 216-room motel available to transient guests on Courtland Street, two blocks from downtown Peachtree Street in Atlanta, Georgia (\href{https://www.law.cornell.edu/supremecourt/text/379/241}{Cornell LII}). The motel was readily accessible to Interstate Highways 75 and 85 and State Highways 23 and 41 (\href{https://supreme.justia.com/cases/federal/us/379/241/}{Justia}). It solicited patronage from outside Georgia through national advertising media, including magazines of national circulation, maintained more than 50 billboards and highway signs within Georgia, accepted convention trade from outside Georgia, and drew approximately 75\% of its registered guests from out of state (\href{https://supreme.justia.com/cases/federal/us/379/241/}{Justia}).

Prior to the passage of the Civil Rights Act of 1964, the motel followed a policy of refusing to rent rooms to Black (``Negro'') guests (\href{https://supreme.justia.com/cases/federal/us/379/241/}{Justia}). After the Act took effect, the motel continued to refuse transient Black guests and alleged that it intended to perpetuate that policy unless restrained (\href{https://www.law.cornell.edu/supremecourt/text/379/241}{Cornell LII}). The motel submitted the case on pleadings, admissions, and stipulated facts, challenging Title II under the Commerce Clause (U.S.\ Const.\ art.\ I, \S~8, cl.\ 3), the Fifth Amendment Due Process and Takings Clauses, and the Thirteenth Amendment's prohibition of involuntary servitude (\href{https://supreme.justia.com/cases/federal/us/379/241/}{Justia}).

\section{Procedural History}
Appellant filed a declaratory judgment action under 28 U.S.C.\ \S\S~2201--2202 in the U.S.\ District Court for the Northern District of Georgia seeking a declaration that Title II was unconstitutional, an injunction against its enforcement, and damages (\href{https://www.law.cornell.edu/supremecourt/text/379/241}{Cornell LII}). The United States counterclaimed for injunctive enforcement under \S~206(a) of the Act and requested a three-judge district court under \S~206(b) and 28 U.S.C.\ \S~2282 (\href{https://supreme.justia.com/cases/federal/us/379/241/}{Justia}).

On July 22, 1964, a three-judge panel (Circuit Judge Elbert Tuttle and District Judges Frank Hooper and Lewis Morgan) issued a \textit{per curiam} decision upholding \S\S~201(a), 201(b)(1), and 201(c)(1) and entering a permanent injunction on the government's counterclaim barring the motel from refusing Black guests on the basis of race or color (\href{https://law.justia.com/cases/federal/district-courts/FSupp/231/393/1444943/}{\textit{Heart of Atlanta Motel, Inc.\ v.\ United States}, 231 F.\ Supp.\ 393 (N.D.\ Ga.\ 1964)}). Justice Hugo L.\ Black, acting in chambers as Circuit Justice, denied the motel's application for a stay of the injunction (85 S.\ Ct.\ 1 (1964)) (\href{https://supreme.justia.com/cases/federal/us/379/241/}{Justia}). On direct appeal pursuant to \S~206(b), the U.S.\ Supreme Court noted probable jurisdiction (379 U.S.\ 803), heard argument on October 5, 1964 (alongside the companion case \textit{Katzenbach v.\ McClung}, 379 U.S.\ 294 (1964)), and affirmed on December 14, 1964 (\href{https://www.oyez.org/cases/1964/515}{Oyez}).

\section{Judicial Votes}
\begin{itemize}
\item \textbf{Judgment:} Unanimous (9--0) affirmance of the three-judge District Court (\href{https://www.oyez.org/cases/1964/515}{Oyez}).
\item \textbf{Opinion of the Court (8 Justices):} Justice Tom C.\ Clark delivered the opinion of the Court, joined by Chief Justice Earl Warren and Justices William O.\ Douglas, John Marshall Harlan II, William J.\ Brennan, Jr., Potter Stewart, Byron R.\ White, and Arthur J.\ Goldberg (\href{https://supreme.justia.com/cases/federal/us/379/241/}{Justia}).
\item \textbf{Concurring Opinions:}
  \begin{itemize}
  \item Justice Hugo L.\ Black concurred separately in the judgment and holding (applying to both \textit{Heart of Atlanta Motel} and \textit{Katzenbach v.\ McClung}) (\href{https://supreme.justia.com/cases/federal/us/379/241/}{Justia}).
  \item Justice William O.\ Douglas joined the Court's opinion and filed a separate concurrence (also applying to \textit{McClung}) (\href{https://supreme.justia.com/cases/federal/us/379/241/}{Justia}).
  \item Justice Arthur J.\ Goldberg joined the Court's opinion and judgment and filed a separate concurrence (also applying to \textit{McClung}) (\href{https://supreme.justia.com/cases/federal/us/379/241/}{Justia}).
  \end{itemize}
\item \textbf{Dissents:} None (\href{https://www.oyez.org/cases/1964/515}{Oyez}).
\end{itemize}

\section{Holding}
\textbf{Judgment:} Affirmed. The permanent injunction restraining the Heart of Atlanta Motel from refusing to accommodate Black guests under Title II of the Civil Rights Act of 1964 was upheld (\href{https://supreme.justia.com/cases/federal/us/379/241/}{Justia}).

\textbf{Rule Announced:} Congress possesses authority under the Commerce Clause (Art.\ I, \S~8, cl.\ 3) to prohibit racial discrimination by places of public accommodation serving interstate travelers whenever (1) Congress has a rational basis for finding that racial discrimination by such establishments affects interstate commerce, and (2) the means chosen to eliminate that obstruction are reasonable and appropriate (\href{https://www.law.cornell.edu/supremecourt/text/379/241}{Cornell LII}). Because the interstate movement of persons is ``commerce'' regardless of whether the travel is commercial in character, Congress may regulate local incidents of interstate travel that exert a substantial and harmful effect on that commerce, notwithstanding that Congress is simultaneously addressing a moral and social wrong (\href{https://supreme.justia.com/cases/federal/us/379/241/}{Justia}). Furthermore, Title II violates neither the Fifth Amendment (due process or takings) nor the Thirteenth Amendment (involuntary servitude) (\href{https://supreme.justia.com/cases/federal/us/379/241/}{Justia}).

\textbf{Holding vs.\ Dicta:} The Court's strict holding is confined to the validity of Title II (\S\S~201(a), 201(b)(1), and 201(c)(1)) under the Commerce Clause as applied to a motel serving interstate travelers, along with the rejection of the motel's Fifth and Thirteenth Amendment defenses (\href{https://www.law.cornell.edu/supremecourt/text/379/241}{Cornell LII}). The Court expressly declined to pass upon whether Congress also had authority to enact Title II under \S~5 of the Fourteenth Amendment or to adjudicate the state-action provisions in \S~201(d) and \S~202 (\href{https://supreme.justia.com/cases/federal/us/379/241/}{Justia}). Likewise, Justice Clark's observation that the motel would probably not suffer long-run economic loss from desegregation is dicta, as the Court immediately added that economic loss ``is of no consequence'' to the constitutional inquiry (\href{https://supreme.justia.com/cases/federal/us/379/241/}{Justia}).

\section{Analysis of Opinions}

\subsection{Majority Opinion (Justice Clark)}
Justice Clark began by distinguishing the \href{https://fedsoc.org/case/the-civil-rights-cases}{\textit{Civil Rights Cases}, 109 U.S.\ 3 (1883)}, which had invalidated the public-accommodations provisions of the Civil Rights Act of 1875 (\href{https://supreme.justia.com/cases/federal/us/379/241/}{Justia}). Unlike the 1875 statute, Title II carefully limited its reach to enterprises directly and substantially related to interstate commerce (or supported by state action); moreover, the 1883 Court had not considered the Commerce Clause, and interstate mobility had expanded dramatically since 1875 (\href{https://www.law.cornell.edu/supremecourt/text/379/241}{Cornell LII}).

Although the final Act omitted formal statutory findings, Justice Clark emphasized that the congressional hearings were ``replete with evidence'' that racial discrimination in lodging inflicted both qualitative and quantitative burdens on Black interstate travelers---forcing them to rely on special travel guidebooks, drive long distances, or forgo interstate travel altogether (\href{https://supreme.justia.com/cases/federal/us/379/241/}{Justia}). Drawing on Chief Justice Marshall's foundational definition of commerce in \textit{Gibbons v.\ Ogden}, 22 U.S.\ (9 Wheat.) 1 (1824), and a long line of precedents establishing that the interstate movement of persons is commerce (including the \textit{Passenger Cases}, \textit{Hoke v.\ United States}, and \textit{Caminetti v.\ United States}), Clark concluded that Congress's plenary commerce power extends to removing local obstructions to interstate travel, even when Congress is also legislating against a moral wrong (\href{https://www.law.cornell.edu/supremecourt/text/379/241}{Cornell LII}).

Finally, the majority dispatched the motel's individual-rights defenses. Under the Fifth Amendment, 32 states already barred public-accommodations discrimination without violating due process, and consequential economic loss from valid commercial regulation is not a compensable taking (\href{https://supreme.justia.com/cases/federal/us/379/241/}{Justia}). Under the Thirteenth Amendment, Title II merely codified the centuries-old common-law innkeeper rule and bore no resemblance to ``African slavery'' (\href{https://supreme.justia.com/cases/federal/us/379/241/}{Justia}).

\subsection{Concurring Opinion (Justice Black)}
Writing separately for both \textit{Heart of Atlanta Motel} and \textit{Katzenbach v.\ McClung}, Justice Black anchored Title II jointly in the Commerce Clause and the Necessary and Proper Clause (Art.\ I, \S~8, cl.\ 18), invoking \textit{McCulloch v.\ Maryland} and the \textit{Shreveport Rate Cases} (\href{https://supreme.justia.com/cases/federal/us/379/241/}{Justia}). Black cautioned that ``every remote, possible, speculative effect on commerce'' cannot erase the line between local and national spheres, and noted that whether an activity sufficiently affects interstate commerce remains ultimately a judicial question (\href{https://www.law.cornell.edu/supremecourt/text/379/241}{Cornell LII}). Applying the aggregation principle of \textit{Wickard v.\ Filburn}, however, he concluded that widespread racial exclusion by motels and restaurants imposed a severe aggregate burden on interstate travel and trade (\href{https://supreme.justia.com/cases/federal/us/379/241/}{Justia}). He rejected the Fifth Amendment claim as ``highly ironical'' and agreed there was no need to decide whether \S~5 of the Fourteenth Amendment independently supported the Act (\href{https://supreme.justia.com/cases/federal/us/379/241/}{Justia}).

\subsection{Concurring Opinion (Justice Douglas)}
While joining the Court's opinion, Justice Douglas expressed reluctance to rest civil rights legislation solely on the Commerce Clause, arguing that the right to be free of racial discrimination ``occupies a more protected position in our constitutional system than does the movement of cattle, fruit, steel and coal across state lines'' (\href{https://supreme.justia.com/cases/federal/us/379/241/}{Justia}). Relying on \S~5 of the Fourteenth Amendment and \textit{Shelley v.\ Kraemer}, 334 U.S.\ 1 (1948), Douglas contended that state enforcement of trespass laws to back up discriminatory exclusions constituted state action (\href{https://www.law.cornell.edu/supremecourt/text/379/241}{Cornell LII}). Grounding Title II in the Fourteenth Amendment, he argued, would have a ``more settling effect,'' eliminating case-by-case litigation over whether a particular inn or restaurant satisfied commerce thresholds (\href{https://supreme.justia.com/cases/federal/us/379/241/}{Justia}).

\subsection{Concurring Opinion (Justice Goldberg)}
Justice Goldberg joined the Court's opinion and judgment but wrote to underscore that the primary purpose of the Civil Rights Act of 1964 ``is the vindication of human dignity, and not mere economics'' (\href{https://supreme.justia.com/cases/federal/us/379/241/}{Justia}). Quoting the Senate Commerce Committee report and his own concurrence in \textit{Bell v.\ Maryland}, 378 U.S.\ 226 (1964), Goldberg maintained that Congress clearly possessed authority under \textit{both} the Commerce Clause and \S~5 of the Fourteenth Amendment to enforce the constitutional right to equal access to public accommodations (\href{https://www.law.cornell.edu/supremecourt/text/379/241}{Cornell LII}).

\subsection{Dissents}
No Justice dissented (\href{https://www.oyez.org/cases/1964/515}{Oyez}).

\section{Examples: Future Applications}
\begin{enumerate}
\item \textbf{Same-Side (Valid under Commerce Clause):} Congress enacts a federal statute prohibiting private short-term rental platforms and commercial vacation-rental hosts with more than five units from discriminating based on disability or familial status, supported by hearings showing that discrimination deters interstate travel. \textit{Analysis:} Like the Heart of Atlanta Motel, commercial lodging offered to transient travelers is an economic activity directly tied to interstate travel; under \textit{Heart of Atlanta Motel}, Congress's moral objective does not negate its Commerce Clause authority (\href{https://supreme.justia.com/cases/federal/us/379/241/}{Justia}).
\item \textbf{Same-Side (Valid under Commerce Clause):} Congress passes a law barring privately owned electric-vehicle charging stations along federal and state highways from refusing service to drivers based on out-of-state vehicle registration or protected characteristics. \textit{Analysis:} Charging stations serving highway motorists are direct instrumentalities and local incidents of interstate travel, squarely within the rational-basis framework of \textit{Heart of Atlanta Motel} (\href{https://www.law.cornell.edu/supremecourt/text/379/241}{Cornell LII}).
\item \textbf{Opposite-Side (Invalid under Commerce Clause):} Congress creates a federal civil remedy for victims of purely local, non-commercial, gender-motivated interpersonal violence, relying on legislative findings that fear of such violence reduces interstate travel and economic productivity. \textit{Analysis:} Under the Rehnquist Court's boundary in \href{https://supreme.justia.com/cases/federal/us/529/598/}{\textit{United States v.\ Morrison}, 529 U.S.\ 598 (2000)}, gender-motivated violent crime is non-economic activity; unlike the commercial motel operations in \textit{Heart of Atlanta Motel}, Congress cannot aggregate attenuated effects of non-economic violent crime to assert general police power.
\item \textbf{Opposite-Side (Invalid under Commerce Clause):} Congress enacts a statute criminalizing the simple possession of a firearm within 1,000 feet of a local public school without any jurisdictional nexus to interstate commerce. \textit{Analysis:} As in \textit{United States v.\ Lopez}, 514 U.S.\ 549 (1995), discussed in the \href{https://constitution.heritage.org/essays/a1-s8-c3-b/}{Heritage Foundation Commerce Clause analysis}, mere possession of a firearm near a school is neither commercial activity nor a local incident of interstate travel, falling outside the outer bounds even of \textit{Heart of Atlanta Motel}.
\item \textbf{Fence-Sitter (Borderline):} Congress amends Title II to eliminate the five-room ``Mrs.\ Murphy'' exemption and require an owner-occupied two-bedroom farmhouse in a remote rural county---which advertises only on a local community bulletin board and hosts almost exclusively in-state relatives and hunters---to accept all transient guests without discrimination. \textit{Analysis:} On one hand, renting rooms for money is commercial activity whose aggregate exclusion across all small boardinghouses could affect interstate travel (echoing Justice Black's aggregation reasoning and \textit{Katzenbach v.\ McClung}). On the other hand, the absence of out-of-state marketing or interstate highway proximity tests Justice Black's warning that ``remote, possible, speculative'' local effects should not obliterate the constitutional distinction between local and national spheres (\href{https://supreme.justia.com/cases/federal/us/379/241/}{Justia}).
\end{enumerate}

\section{Critique}
\textbf{Originalist and Textualist Perspectives:} Originalist and libertarian scholars such as Randy Barnett and Richard Epstein, alongside Justice Clarence Thomas in later Rehnquist Court opinions, have questioned the expansive ``substantial effects'' doctrine that undergirded \textit{Heart of Atlanta Motel} (\href{https://digitalcommons.law.byu.edu/cgi/viewcontent.cgi?article=3408&context=lawreview}{BYU Law Review}; \href{https://dspace.loyno.edu/jspui/bitstream/123456789/131/1/Article\%204\%20-\%20McGoldrick_PRINT.pdf}{Loyola Law Review}). From a strict original-meaning view of ``Commerce,'' the term referred primarily to trade and the cross-border movement of goods rather than intrastate service businesses, raising the concern that deferential rational-basis review converts the Commerce Clause into a plenary federal police power (\href{https://constitution.heritage.org/essays/a1-s8-c3-b/}{Heritage Foundation}). Richard Epstein has further argued that broad public-accommodations mandates under Title II override freedom of association and property rights beyond the narrow common-law monopoly rationale for innkeepers (\href{https://chicagounbound.uchicago.edu/cgi/viewcontent.cgi?article=8104&context=journal_articles}{Epstein, 66 Stan.\ L.\ Rev.\ 1241 (2014)}).

\textbf{Progressive and Doctrinal Perspectives:} Many progressive constitutional scholars applaud the decision's result while echoing Justices Douglas and Goldberg in criticizing its constitutional foundation. By resting Title II solely on the Commerce Clause and leaving the \href{https://fedsoc.org/case/the-civil-rights-cases}{\textit{Civil Rights Cases} (1883)} intact, the Court framed racial equality in commercial terms rather than vindicating human dignity under \S~5 of the Fourteenth Amendment (\href{https://scholarship.law.tamu.edu/cgi/viewcontent.cgi?article=1359&context=lawreview}{Texas A\&M Law Review}). That doctrinal choice carried long-term consequences during the Rehnquist Court era: when the Court later tightened both Commerce Clause limits and \S~5 enforcement power in \href{https://supreme.justia.com/cases/federal/us/529/598/}{\textit{United States v.\ Morrison} (2000)} and \textit{City of Boerne v.\ Flores} (1997), Congress lacked an entrenched Fourteenth Amendment precedent for regulating private discrimination (\href{https://scholarlycommons.law.northwestern.edu/cgi/viewcontent.cgi?article=1349&context=nulr}{Northwestern Law Review}).

\textbf{Underlying Competing Values:} The case sits at the intersection of four enduring constitutional tensions: (1) enumerated federal powers and federalism versus the need for uniform national civil rights enforcement; (2) private property and freedom of association versus equal citizenship and human dignity in the public marketplace; (3) judicial deference to congressional fact-finding versus judicial enforcement of structural limits on Article I power; and (4) pragmatic doctrinal consensus (avoiding overruling the \textit{Civil Rights Cases}) versus doctrinal candor about the true moral basis of civil rights legislation (\href{https://fedsoc.org/case/heart-of-atlanta-motel-inc-v-united-states}{Federalist Society}).

\section{Key Quotations}
\begin{quote}
``In short, the determinative test of the exercise of power by the Congress under the Commerce Clause is simply whether the activity sought to be regulated is `commerce which concerns more States than one' and has a real and substantial relation to the national interest.''
\end{quote}
\noindent --- Justice Tom C.\ Clark, Opinion of the Court, 379 U.S.\ at 255 (\href{https://supreme.justia.com/cases/federal/us/379/241/}{Justia}).

\begin{quote}
``It was this burden which empowered Congress to enact appropriate legislation, and, given this basis for the exercise of its power, Congress was not restricted by the fact that the particular obstruction to interstate commerce with which it was dealing was also deemed a moral and social wrong.''
\end{quote}
\noindent --- Justice Tom C.\ Clark, Opinion of the Court, 379 U.S.\ at 257 (\href{https://supreme.justia.com/cases/federal/us/379/241/}{Justia}).

\begin{quote}
``Thus, the power of Congress to promote interstate commerce also includes the power to regulate the local incidents thereof, including local activities in both the States of origin and destination, which might have a substantial and harmful effect upon that commerce.''
\end{quote}
\noindent --- Justice Tom C.\ Clark, Opinion of the Court, 379 U.S.\ at 258 (\href{https://supreme.justia.com/cases/federal/us/379/241/}{Justia}).

\begin{quote}
``A decision based on the Fourteenth Amendment would have a more settling effect, making unnecessary litigation over whether a particular restaurant or inn is within the commerce definitions of the Act or whether a particular customer is an interstate traveler.''
\end{quote}
\noindent --- Justice William O.\ Douglas, concurring, 379 U.S.\ at 280 (\href{https://supreme.justia.com/cases/federal/us/379/241/}{Justia}).

\begin{quote}
``The primary purpose of the Civil Rights Act of 1964, however, as the Court recognizes, and as I would underscore, is the vindication of human dignity, and not mere economics.''
\end{quote}
\noindent --- Justice Arthur J.\ Goldberg, concurring, 379 U.S.\ at 291 (\href{https://www.law.cornell.edu/supremecourt/text/379/241}{Cornell LII}).

\end{document}
